\chapter{Primi programmi in C}
\label{chap:primi-programmi}

Questo capitolo è una panoramica del C attraverso una serie di programmi di
esempio, ciascuno spiegato elemento per elemento. Si introducono
brevemente i costrutti principali (l'output, le variabili, l'input, le
scelte, i cicli, le funzioni) per avere subito un'idea di come è fatto un
programma vero; ogni argomento verrà poi approfondito nei capitoli
successivi. Il consiglio è uno solo: \textbf{scrivere ed eseguire ogni
esempio}, poi modificarlo e osservare che cosa cambia.

\section{Il primo programma, riga per riga}

\subsection{Scrivere, compilare, eseguire}

Un programma deve comunicare per essere utile: produce sempre un output di
qualche tipo. Il primo esempio stampa una frase sullo schermo. Lo scriviamo
nel file \code{sea.c}, lo compiliamo e lo eseguiamo, come abbiamo imparato
nel capitolo~\ref{chap:ambiente}:

\begin{affianco}[0.5]
\begin{lstlisting}[style=c, name=sea.c]
#include <stdio.h>

int main(void)
{
    printf("from sea to shining C\n");
    return 0;
}
\end{lstlisting}
\risultato
\begin{console}
$ gcc sea.c
$ ./a.out
from sea to shining C
\end{console}
\wtnota{Se non ci sono errori, gcc crea il file eseguibile \code{a.out}.}
\end{affianco}

Sono sette righe, ma ognuna ha un ruolo preciso. La
figura~\ref{fig:anatomia} le numera, e l'elenco che la segue le spiega.

\begin{figure}[H]
  \centering
  \begin{tikzpicture}[font=\FiraCodeNL\small,
      riga/.style={anchor=base west, inner sep=0pt},
      num/.style={circle, fill=accent!75!black, text=white, font=\Lato\scriptsize\bfseries,
        inner sep=0pt, minimum size=0.42cm}]
    \fill[vscbg, rounded corners=4pt] (-0.3,0.55) rectangle (9.2,-3.4);
    \node[riga] (r1) at (0,0) {\textcolor{vscpurple}{\#include} \textcolor{vscstring}{<stdio.h>}};
    \node[riga] (r2) at (0,-0.6) {\textcolor{vscblue}{int} \textcolor{vscfunc}{main}\textcolor{vscfg}{(}\textcolor{vscblue}{void}\textcolor{vscfg}{)}};
    \node[riga, text=vscfg] (r3) at (0,-1.2) {\{};
    \node[riga] (r4) at (0,-1.8) {\textcolor{vscfg}{\ \ \ \ }\textcolor{vscfunc}{printf}\textcolor{vscfg}{(}\textcolor{vscstring}{"from sea to shining C\textbackslash n"}\textcolor{vscfg}{);}};
    \node[riga] (r5) at (0,-2.4) {\textcolor{vscfg}{\ \ \ \ }\textcolor{vscpurple}{return} \textcolor{vscfg}{0;}};
    \node[riga, text=vscfg] (r6) at (0,-3.0) {\}};
    \foreach \r/\n in {r1/1, r2/2, r3/3, r4/4, r5/5, r6/3}
      \node[num] at ([yshift=3.5pt]8.75,0 |- \r.base) {\n};
  \end{tikzpicture}
  \caption{Anatomia del primo programma: i numeri rimandano all'elenco
           seguente.}
  \label{fig:anatomia}
\end{figure}

\begin{passi}
  \item \textbf{Direttiva per il preprocessore}: copia qui il contenuto del
        file di intestazione \code{stdio.h}, che descrive \code{printf}.
  \item \textbf{Testa della funzione} \code{main}, da cui comincia
        l'esecuzione: restituisce un \code{int} e non riceve parametri
        (\code{void}).
  \item Le \textbf{graffe} delimitano il \textbf{corpo} della funzione.
  \item Un'\textbf{istruzione}: chiama \code{printf} con un argomento, una
        \textbf{stringa costante}; \code{\textbackslash n} va a capo; il
        \code{;} chiude l'istruzione.
  \item Restituisce 0 al sistema operativo: “tutto bene”.
\end{passi}

\subsection{Il preprocessore e \code{\#include}}

Prima della compilazione vera e propria viene eseguito, automaticamente, il
\term{preprocessore} (capitolo~\ref{chap:dal-problema}). Le righe che
cominciano con \code{\#} sono \textbf{direttive} per il preprocessore. La
direttiva \code{\#include <stdio.h>} chiede di includere in quel punto una
copia del file di intestazione (\textit{header file}) \code{stdio.h}, che
contiene le informazioni necessarie per usare le funzioni di input e output,
come \code{printf}.

\warning{\code{\#include} non include la libreria}{
  A lezione il docente è categorico: guai a dire che \code{\#include}
  “include la libreria”. Include un \textbf{file di testo}, l'header, che
  \textbf{descrive} le funzioni (come si chiamano, che parametri vogliono,
  che cosa restituiscono). Il \textbf{codice} delle funzioni sta nella
  libreria, che viene collegata dal \textbf{linker}, dopo. Sono due cose
  diverse, che intervengono in due fasi diverse: lo vedremo in modo molto
  concreto con la funzione \code{pow} più avanti in questo capitolo.}

\subsection{La funzione \code{main}}

\code{int main(void)} è la prima riga della definizione di una
\term{funzione} chiamata \code{main}. Le parentesi dopo il nome dicono che
\code{main} è una funzione; \code{int} e \code{void} sono parole speciali
che danno informazioni al compilatore: \code{void} fra le parentesi
significa che la funzione \textbf{non ha parametri}, \code{int} davanti al
nome significa che \textbf{restituisce un intero}.

Ogni programma C ha \textbf{una} funzione \code{main}, e l'esecuzione
comincia da lei. Il corpo della funzione è racchiuso fra le graffe
\code{\{} e \code{\}}, che in C servono in generale a raggruppare
istruzioni.

\subsection{\code{printf}, le stringhe e il punto e virgola}

\code{printf} è una funzione che stampa sullo schermo (più precisamente
sullo \textit{standard output}, come vedremo). Non fa parte del
linguaggio: appartiene alla \term{libreria standard}, un insieme di
librerie che forniscono al programmatore funzioni pronte; gli header file
permettono di usarle, e \code{stdio.h} fornisce le informazioni per
\code{printf}.

Il testo fra virgolette, \code{"from sea to shining C\textbackslash n"}, è
una \term{stringa costante}: è l'\textbf{argomento} di \code{printf} e
stabilisce che cosa viene stampato. Al suo interno, \code{\textbackslash n}
(una barra rovesciata seguita da \code{n}) è un \textbf{singolo carattere},
il \textit{newline}: fa avanzare il cursore all'inizio della riga
successiva. L'intera riga è un'\term{istruzione}, e in C le istruzioni
\textbf{terminano con un punto e virgola}.

\subsection{\code{return} e il valore restituito al sistema}

\code{return 0;} termina la funzione \code{main} e restituisce il valore 0
al sistema operativo. Per convenzione 0 significa “tutto è andato bene”,
un valore diverso da 0 segnala un problema. Il valore non si vede
nell'output, ma la shell lo conserva, e lo si può leggere con il comando
\code{echo \$?} subito dopo l'esecuzione:

\begin{console}
$ ./a.out
from sea to shining C
$ echo $?
0
\end{console}

\code{return} \textbf{termina} la funzione: ciò che viene dopo non è
eseguito. Il file \code{sea.c} distribuito con le lezioni ne dà una prova,
perché dopo il \code{return} contiene un'altra \code{printf}:

\begin{affianco}[0.5]
\begin{lstlisting}[style=c, name=sea.c]
#include <stdio.h>

int main()
{
    printf("from sea to shining C\n");

    return 0;
    printf("Pippo\n");
}
\end{lstlisting}
\risultato
\begin{console}
$ gcc -o sea sea.c
$ ./sea
from sea to shining C
\end{console}
\wtnota{“Pippo” non viene mai stampato: l'esecuzione esce da \code{main}
al \code{return}.}
\end{affianco}

(Qui \code{main} è scritto con le parentesi vuote, \code{main()}: per una
definizione di funzione è equivalente a \code{main(void)}; in questa
dispensa useremo la forma con \code{void}, più esplicita.)

\section{L'output con \code{printf}}

\subsection{Il cursore}

\code{printf} scrive sullo schermo da sinistra a destra e dall'alto verso
il basso, come si scrive su un quaderno; va a capo \textbf{solo} quando
incontra il carattere \code{\textbackslash n}. Ogni chiamata di
\code{printf} comincia da dove la precedente ha lasciato il cursore. Per
questo il programma seguente, diverso dal primo, stampa \textbf{lo stesso}
output:

\begin{affianco}[0.5][c]
\begin{lstlisting}[style=c]
printf("from sea to ");
printf("shining ");
printf("C\n");
\end{lstlisting}
\risultato
\begin{console}
from sea to shining C
\end{console}
\end{affianco}

Se vogliamo scrivere su tre righe dobbiamo inserire dei \code{\textbackslash
n}; e se ne possono mettere più di uno nella stessa stringa:

\begin{affianco}[0.5][c]
\begin{lstlisting}[style=c]
printf("from sea to\n");
printf("shining\nC\n");
\end{lstlisting}
\risultato
\begin{console}
from sea to
shining
C
\end{console}
\end{affianco}

\subsection{Disegnare con i caratteri}

Per essere leggibile, l'output va spaziato opportunamente sullo schermo.
Il programma seguente stampa la stessa frase in una cornice di asterischi,
preceduta e seguita da dieci righe vuote (ciascun \code{\textbackslash n}
iniziale produce una riga vuota):

\begin{affianco}[0.58][c]
\begin{lstlisting}[style=c, name=sea_box.c]
#include <stdio.h>

int main(void)
{
    printf("\n\n\n\n\n\n\n\n\n\n");
    printf("      ***********************\n");
    printf("      *  from sea           *\n");
    printf("      *  to shining C       *\n");
    printf("      ***********************\n");
    printf("\n\n\n\n\n\n\n\n\n\n");
    return 0;
}
\end{lstlisting}
\risultato
\begin{console}
$ ./sea_box


      ***********************
      *  from sea           *
      *  to shining C       *
      ***********************


$
\end{console}
\wtnota{Qui sono mostrate solo due delle dieci righe vuote sopra e sotto.}
\end{affianco}

\section{Variabili, espressioni e assegnazione}

\subsection{Il problema della maratona}

Il secondo esempio converte una distanza: la lunghezza di una maratona è di
26 miglia e 385 iarde, e vogliamo sapere quanti chilometri sono. Servono due
regole: in un miglio ci sono 1760 iarde, e un miglio è lungo 1,609 km.

\begin{affianco}[0.56]
\begin{lstlisting}[style=c, name=marathon.c]
/* The distance of a marathon in kilometers. */

#include <stdio.h>

int main(void)
{
    int     miles, yards;
    float   kilometers;

    miles = 26;
    yards = 385;
    kilometers = 1.609 * (miles + yards / 1760.0);
    printf("\nA marathon is %f kilometers.\n\n", kilometers);
    return 0;
}
\end{lstlisting}
\risultato
\begin{console}
$ gcc -o marathon marathon.c
$ ./marathon

A marathon is 42.185970 kilometers.

\end{console}
\wtnota{L'opzione \code{-o} sceglie il nome del file eseguibile.}
\end{affianco}

\subsection{I commenti}

Tutto ciò che è compreso fra \code{/*} e \code{*/} è un
\term{commento}: viene ignorato dal compilatore (anzi, eliminato già dal
preprocessore) e serve solo a chi legge. Si possono anche scrivere commenti
su una sola riga, che cominciano con \code{//} e arrivano fino alla fine
della riga. Un commento all'inizio del file spiega lo scopo del programma;
vi si possono aggiungere informazioni sugli autori e sulle versioni.

\subsection{Le dichiarazioni}

\begin{lstlisting}[style=c]
int     miles, yards;
float   kilometers;
\end{lstlisting}

Queste righe sono \term{dichiarazioni}. \code{int} e \code{float} sono
\textbf{parole chiave} del C che indicano due \textbf{tipi di dato
fondamentali}: la prima riga dice al compilatore che \code{miles} e
\code{yards} sono \term{variabili} di tipo intero, che assumono valori
interi; la seconda che \code{kilometers} è una variabile di tipo
\code{float}, che assume valori reali. Anche le dichiarazioni terminano con
un punto e virgola.

Dichiarare una variabile significa chiedere al compilatore di riservarle
dello spazio in memoria e di ricordarne il tipo. In C \textbf{ogni
variabile va dichiarata prima di essere usata}.

\subsection{L'assegnazione}

\code{miles = 26;} è un'\term{istruzione di assegnazione}: il simbolo
\code{=} è l'\textbf{operatore di assegnazione}, \code{26} è una costante
intera, e l'effetto è che il valore 26 viene scritto nella variabile
\code{miles} (figura~\ref{fig:maratona-memoria}). Allo stesso modo 385 viene
assegnato a \code{yards}.

\warning{\code{=} non è l'uguale della matematica}{
  In matematica $x = x + 1$ è un'equazione impossibile. In C
  \code{x = x + 1;} è perfettamente sensato e significa: \textbf{calcola} il
  valore dell'espressione a destra (il valore attuale di \code{x} più 1) e
  \textbf{scrivilo} nella variabile a sinistra. L'assegnazione va letta da
  destra a sinistra: “\code{x} riceve\dots”.}

\begin{figure}[H]
  \centering
  \begin{tikzpicture}[font=\Lato\scriptsize]
    \foreach \titolo/\m/\y/\km [count=\i from 0] in {
        {dopo le dichiarazioni}/\mvuota/\mvuota/\mvuota,
        {dopo \code{miles = 26;}}/26/\mvuota/\mvuota,
        {dopo \code{yards = 385;}}/26/385/\mvuota,
        {dopo il calcolo}/26/385/42.18597} {
      \begin{scope}[xshift=\i*4.4cm]
        \node[font=\Lato\scriptsize\bfseries] at (0.3,0.6) {\titolo};
        \variabile{m\i}{(0.6,0)}{miles}{\m}
        \variabile{y\i}{(0.6,-0.65)}{yards}{\y}
        \variabile[minimum width=1.4cm]{k\i}{(0.6,-1.3)}{kilometers}{\km}
      \end{scope}
    }
  \end{tikzpicture}
  \caption{Le variabili della maratona, istruzione dopo istruzione. Appena
           dichiarata, una variabile contiene un valore a caso (“?”):
           finché non le si assegna qualcosa, non va usata.}
  \label{fig:maratona-memoria}
\end{figure}

\subsection{Come si calcola un'espressione}

\code{kilometers = 1.609 * (miles + yards / 1760.0);} assegna un valore
calcolato da un'\term{espressione}. Le regole sono quelle dell'aritmetica:
si calcola prima ciò che sta fra parentesi, e dentro le parentesi la
divisione ha la precedenza sulla somma. Quindi prima \code{yards / 1760.0},
poi il risultato viene sommato a \code{miles}, e infine moltiplicato per
\code{1.609} (figura~\ref{fig:albero-maratona}).

\begin{figure}[H]
  \centering
  \begin{tikzpicture}[font=\FiraCodeNL\small,
      op/.style={circle, draw=accent!65!black, fill=accent!8!white, line width=0.8pt,
        minimum size=0.62cm, inner sep=0pt, font=\FiraCodeNL\small\bfseries},
      fo/.style={wtnodeplain, font=\FiraCodeNL\scriptsize, minimum width=1.2cm},
      val/.style={font=\Lato\scriptsize\bfseries, text=green-gradient-6!50!black, anchor=west},
      ord/.style={circle, fill=primary!75!black, text=white, font=\Lato\tiny\bfseries,
        inner sep=0pt, minimum size=0.36cm}]
    \node[op] (mul) at (0,0) {*};
    \node[fo] (k) at (-1.6,-1.1) {1.609};
    \node[op] (add) at (1.6,-1.1) {+};
    \node[fo] (m) at (0.4,-2.2) {miles\\[-1pt]{\color{neutral!50!black}(26)}};
    \node[op] (div) at (2.8,-2.2) {/};
    \node[fo] (y) at (1.8,-3.35) {yards\\[-1pt]{\color{neutral!50!black}(385)}};
    \node[fo] (c) at (3.8,-3.35) {1760.0};
    \foreach \a/\b in {mul/k, mul/add, add/m, add/div, div/y, div/c} \draw[fsramo] (\a) -- (\b);
    \node[val] at ([xshift=4pt]div.east) {= 0.21875};
    \node[val] at ([xshift=4pt]add.east) {= 26.21875};
    \node[val] at ([xshift=4pt]mul.east) {= 42.18597};
    \node[ord] at ([shift={(-0.42,0.32)}]div.center) {1};
    \node[ord] at ([shift={(-0.42,0.32)}]add.center) {2};
    \node[ord] at ([shift={(-0.42,0.32)}]mul.center) {3};
  \end{tikzpicture}
  \caption{L'albero dell'espressione: si calcola dal basso verso l'alto,
           nell'ordine indicato dai numeri.}
  \label{fig:albero-maratona}
\end{figure}

\code{1.609} e \code{1760.0} sono \textbf{costanti reali} (in C, di tipo
\code{double}): per questo \code{yards / 1760.0} è una divisione fra reali,
con risultato reale. Infine \code{printf("\dots\%f\dots", kilometers)}
stampa il valore: la stringa di controllo contiene una
\term{specifica di conversione}, \code{\%f}, che dice “qui va stampato un
numero reale”; la variabile corrispondente, il secondo argomento, deve
essere di quel tipo.

\subsection{La trappola della divisione intera}

Basta togliere il \code{.0} alla costante per cambiare completamente il
risultato:

\begin{affianco}[0.6][c]
\begin{lstlisting}[style=c]
kilometers = 1.609 * (miles + yards / 1760);
\end{lstlisting}
\risultato
\begin{console}
A marathon is 41.834000 kilometers.
\end{console}
\end{affianco}

Ora \code{1760} è una costante \textbf{intera}, e \code{yards / 1760} è
una divisione fra interi: in C il risultato di una divisione fra interi è
un intero, e il resto si perde. 385 diviso 1760 fa 0 con resto 385, quindi
\code{yards / 1760} vale \textbf{0} e i 385 metri e rotti\dots\ spariscono:
$1{,}609 \cdot 26 = 41{,}834$. Il compilatore non segnala niente, perché
non c'è nessun errore di linguaggio: c'è un errore di ragionamento.

\section{Parole chiave e nomi}

Parole come \code{int} e \code{float} sono \term{parole chiave}
(\textit{keywords}): sono \textbf{riservate} e non possono essere usate come
nomi di variabili. Una variabile non può chiamarsi \code{int}, né
\code{return}. Per lo stesso motivo di chiarezza, non si usano come nomi di
variabili i nomi delle funzioni: niente variabili chiamate \code{printf} o
\code{main}. L'elenco completo delle parole chiave e le regole per i nomi
sono nel capitolo~\ref{chap:sintassi}.

\section{Il preprocessore: \code{\#define} e \code{\#include}}

\subsection{Le costanti simboliche}

Il preprocessore scandisce tutto il file e individua le direttive, cioè le
righe che cominciano con \code{\#}. Due delle più usate sono
\code{\#define} e \code{\#include}. Consideriamo le righe

\begin{lstlisting}[style=c]
#define LIMIT 100
#define PI    3.14159
\end{lstlisting}

Il preprocessore sostituirà nel file sorgente tutte le occorrenze di
\code{LIMIT} con \code{100} e tutte quelle di \code{PI} con
\code{3.14159}, tranne quelle che compaiono fra virgolette (nelle stringhe)
e nei commenti. Un \code{\#define} può stare ovunque nel file e ha effetto
solo sulle righe successive; di solito si mette all'inizio, e per
convenzione il nome si scrive in maiuscolo.

Le \term{costanti simboliche} rendono il programma più \textbf{leggibile}
(\code{LIMIT} dice che cosa rappresenta, \code{100} no) e più facile e
sicuro da \textbf{modificare}. Se la velocità della luce, circa 299\,792,458
km/s, compare in cento punti di un programma, scrivendo

\begin{lstlisting}[style=c]
#define C 299792.458   /* velocità della luce in km/s */
\end{lstlisting}

e usando \code{C} ovunque, una stima più accurata si inserisce cambiando
una sola riga, senza rischiare di dimenticare qualche occorrenza.

\subsection{I file di intestazione}

Una riga come

\begin{lstlisting}[style=c]
#include "my_file.h"
\end{lstlisting}

copia il contenuto di \code{my\_file.h} in quel punto del sorgente. Anche
\code{\#include} può stare ovunque, ma tipicamente si mette all'inizio. I
file che si includono sono gli \term{header file} e per convenzione hanno
estensione \code{.h}. Il sistema C ne fornisce alcuni, come
\code{stdio.h}, \code{stdlib.h}, \code{string.h}, \code{math.h}. La
differenza fra le due forme:

\begin{itemize}
  \item \code{\#include <stdio.h>}, con le parentesi angolari, cerca il file
        fra gli header del sistema;
  \item \code{\#include "my\_file.h"}, con le virgolette, lo cerca prima
        nella cartella del sorgente: si usa per i nostri header.
\end{itemize}

Ogni volta che si usa una funzione di input/output come \code{printf} si
deve includere \code{stdio.h}, che ne contiene il \textbf{prototipo}: la
descrizione che serve al compilatore per controllare che la chiamata sia
corretta. L'header serve alla \textbf{compilazione}; il codice della
funzione arriva dalla libreria, durante il \textbf{linking}.

\subsection{Header e libreria: la prova}

La funzione \code{pow(x, y)} calcola $x^y$; il suo prototipo sta in
\code{math.h}, il suo codice nella \textbf{libreria matematica}, che su
Linux è separata dalla libreria standard e va chiesta esplicitamente al
linker con l'opzione \code{-lm} (\textit{library m}):

\begin{affianco}[0.48]
\begin{lstlisting}[style=c, name=potenza.c]
#include <stdio.h>
#include <math.h>

int main(void)
{
    double base, esponente;

    printf("Base ed esponente: ");
    scanf("%lf%lf", &base, &esponente);
    printf("%g elevato a %g fa %g\n",
           base, esponente, pow(base, esponente));
    return 0;
}
\end{lstlisting}
\risultato
\begin{console}
$ gcc -Wall -o potenza potenza.c
/usr/bin/ld: /tmp/ccVCbC9C.o: in function `main':
potenza.c:(.text+0x61): undefined reference to `pow'
collect2: error: ld returned 1 exit status
$ gcc -Wall -o potenza potenza.c -lm
$ ./potenza
Base ed esponente: <<2 10>>
2 elevato a 10 fa 1024
\end{console}
\end{affianco}

Il primo tentativo fallisce, ma si guardi \textbf{chi} protesta:
\code{ld}, il linker, non il compilatore. Grazie a \code{math.h} il
compilatore sapeva tutto ciò che gli serviva e ha prodotto il file oggetto;
è il linker che non ha trovato il \textbf{codice} di \code{pow}
(\textit{undefined reference}). Con \code{-lm} gli si indica la libreria, e
tutto funziona. (Su macOS la libreria matematica è già compresa in quella
di sistema e \code{-lm} non serve, ma non fa danni.)

\section{Un esempio con le costanti: Pacific Sea}

\subsection{Il programma}

\textit{Pacific Sea} è il nome dato alla parte di oceano Pacifico visibile
dal campus di Santa Cruz, sulla baia di Monterey. Il programma converte la
sua area da chilometri quadrati in acri, miglia quadrate, piedi quadrati e
pollici quadrati. Le costanti stanno in un header scritto da noi:

\begin{lstlisting}[style=h, name=pacific_sea.h]
#include <stdio.h>

#define   AREA                        2337
#define   SQ_MILES_PER_SQ_KILOMETER   0.3861021585424458
#define   SQ_FEET_PER_SQ_MILE         (5280 * 5280)
#define   SQ_INCHES_PER_SQ_FOOT       144
#define   ACRES_PER_SQ_MILE           640
\end{lstlisting}

\begin{lstlisting}[style=c, name=pacific_sea.c]
/* Measuring the Pacific Sea. */

#include "pacific_sea.h"

int main(void)
{
    const int   pacific_sea = AREA;     /* in sq kilometers */
    double      acres, sq_miles, sq_feet, sq_inches;

    printf("\nThe Pacific Sea covers an area");
    printf(" of %d square kilometers.\n", pacific_sea);
    sq_miles = SQ_MILES_PER_SQ_KILOMETER * pacific_sea;
    sq_feet = SQ_FEET_PER_SQ_MILE * sq_miles;
    sq_inches = SQ_INCHES_PER_SQ_FOOT * sq_feet;
    acres = ACRES_PER_SQ_MILE * sq_miles;
    printf("In other units of measure this is:\n\n");
    printf("%22.7e acres\n", acres);
    printf("%22.7e square miles\n", sq_miles);
    printf("%22.7e square feet\n", sq_feet);
    printf("%22.7e square inches\n\n", sq_inches);
    return 0;
}
\end{lstlisting}

\begin{console}
$ ./pacific_sea

The Pacific Sea covers an area of 2337 square kilometers.
In other units of measure this is:

         5.7748528e+05 acres
         9.0232074e+02 square miles
         2.5155259e+10 square feet
         3.6223572e+12 square inches

\end{console}

\subsection{Gli \code{\#include} a catena}

La prima riga include una copia di \code{pacific\_sea.h}, e siccome questo
contiene a sua volta \code{\#include <stdio.h>}, il preprocessore include
anche una copia di \code{stdio.h}, che serve per \code{printf}
(figura~\ref{fig:include}). Il risultato è un unico testo, che è quello che
il compilatore riceve.

\begin{figure}[H]
  \centering
  \begin{tikzpicture}[font=\Lato\scriptsize,
      f/.style={wtnodeplain, minimum width=2.9cm, minimum height=1.2cm, align=center}]
    \node[f] (a) at (0,0) {{\FiraCodeNL pacific\_sea.c}\\{\FiraCodeNL\tiny \#include "pacific\_sea.h"}\\il programma};
    \node[f] (b) at (4.6,0) {{\FiraCodeNL pacific\_sea.h}\\{\FiraCodeNL\tiny \#include <stdio.h>}\\le costanti};
    \node[f] (c) at (9.2,0) {{\FiraCodeNL stdio.h}\\{\tiny (header di sistema)}\\i prototipi di I/O};
    \node[wtnode, minimum width=3.0cm, minimum height=1.2cm] (d) at (13.8,0) {un unico testo\\per il compilatore};
    \draw[wtarrowacc] (b) -- node[wtlab, above=3pt, fill=none] {copiato in} (a);
    \draw[wtarrowacc] (c) -- node[wtlab, above=3pt, fill=none] {copiato in} (b);
    \draw[wtarrow] (a.south) -- ++(0,-0.45) -| node[wtlab, pos=0.25] {dopo il preprocessore} (d.south);
  \end{tikzpicture}
  \caption{Il preprocessore risolve gli \code{\#include} uno dentro l'altro.}
  \label{fig:include}
\end{figure}

\subsection{Le costanti dell'header}

Le costanti simboliche dell'header (una reale, le altre intere) rendono il
codice più chiaro: se ci fosse una nuova stima dell'area basterebbe
cambiare \code{AREA}, e le altre danno un nome significativo ai numeri. Su
\code{SQ\_FEET\_PER\_SQ\_MILE} vale la pena fermarsi:

\begin{itemize}
  \item scrivere \code{(5280 * 5280)} invece di \code{27878400} è più
        chiaro: dice che in un miglio ci sono 5280 piedi, e che un miglio
        quadrato ne contiene il quadrato. Non si perde efficienza: il
        compilatore calcola da sé le espressioni fatte solo di costanti;
  \item le parentesi non sono sintatticamente necessarie, ma sono un'ottima
        abitudine. Il preprocessore fa una sostituzione di
        \textbf{testo}, senza capire nulla: se fosse
        \code{\#define SQ\_FEET\_PER\_SQ\_MILE 5280 * 5280}, l'espressione
        \code{x / SQ\_FEET\_PER\_SQ\_MILE} diventerebbe
        \code{x / 5280 * 5280}, cioè (divisione e moltiplicazione si
        eseguono da sinistra a destra) $(x / 5280) \cdot 5280$, che vale di
        nuovo circa $x$. Con le parentesi diventa
        \code{x / (5280 * 5280)}, come si voleva.
\end{itemize}

\subsection{\code{const}}

\code{const int pacific\_sea = AREA;} è la dichiarazione di una variabile
\code{pacific\_sea} di tipo \code{int}, \textbf{inizializzata} (cioè con un
valore assegnato già nella dichiarazione) a 2337: prima della compilazione
il preprocessore ha sostituito \code{AREA} con il suo valore. La parola
chiave \code{const} qualifica il tipo: la variabile riceve il suo valore
all'inizializzazione, ma \textbf{non può essere modificata}. Il compilatore
lo fa rispettare:

\begin{affianco}[0.5]
\begin{lstlisting}[style=c, name=costante.c]
#include <stdio.h>

int main(void)
{
    const int pacific_sea = 2337;

    pacific_sea = 3000;
    printf("%d\n", pacific_sea);
    return 0;
}
\end{lstlisting}
\risultato
\begin{console}
$ gcc -Wall -o costante costante.c
costante.c: In function (*@‘@*)main(*@’@*):
costante.c:7:17: error: assignment of read-only variable (*@‘@*)pacific_sea(*@’@*)
    7 |     pacific_sea = 3000;
      |                 ^
\end{console}
\end{affianco}

\subsection{I tipi reali e la notazione scientifica}

\code{acres}, \code{sq\_miles} e gli altri sono dichiarati \code{double}. In
C i tipi reali (detti \textit{floating point}, “a virgola mobile”) sono
tre: \code{float} (tipicamente 6 cifre significative), \code{double}
(tipicamente 15) e \code{long double} (almeno quanto \code{double}; sui PC
di solito 18). I valori dipendono dall'implementazione; il
capitolo~\ref{chap:tipi} li approfondisce.

L'istruzione \code{printf("\%22.7e acres\textbackslash n", acres);} stampa la
riga \code{5.7748528e+05 acres}. Il numero è in \textbf{notazione
scientifica}: \code{e+05} significa “per $10^5$”, quindi il valore è
$5{,}7748528 \cdot 10^5 = 577\,485{,}28$ acri. La conversione \code{\%e}
chiede proprio questo formato (detto \textit{e-format}), e la forma
\code{\%$m$.$n$e}, con $m$ e $n$ interi positivi, indica che il numero
occupa in tutto $m$ caratteri (allineato a destra) con $n$ cifre dopo il
punto: qui 22 caratteri e 7 cifre decimali. Per questo nell'output i numeri
sono preceduti da spazi e allineati in colonna.

\section{\code{printf} e \code{scanf}}

\subsection{Funzioni di input e output formattato}

\code{printf} si usa per l'output, \code{scanf} per l'input; la \code{f}
finale sta per \textit{formatted} (formattato). Non fanno parte del
linguaggio, ma della libreria standard, che viene collegata implicitamente
alla fine della compilazione. Il programmatore deve comunque dichiarare al
compilatore le funzioni che usa: il meccanismo dell'ANSI C è il
\textbf{prototipo}, e i prototipi delle funzioni della libreria standard
sono negli header standard (per \code{printf} e \code{scanf},
\code{stdio.h}).

Entrambe accettano una lista di argomenti:

\begin{itemize}
  \item la \term{stringa di controllo}, una stringa che contiene le
        \textbf{specifiche di conversione}: ciascuna è il carattere \code{\%}
        seguito da un \textbf{carattere di conversione} (per esempio
        \code{\%d} per un intero);
  \item gli altri argomenti, \textbf{uno per ogni conversione} della stringa
        di controllo, nello stesso ordine.
\end{itemize}

\subsection{Le conversioni di \code{printf}}

Per stampare le prime tre lettere dell'alfabeto ci sono molti modi: senza
conversioni, con tre conversioni per caratteri, con una conversione per
stringhe. Si può anche decidere quanto spazio dare a ciascuna conversione:
\code{\%3c} significa “un carattere, in un campo largo 3”, allineato a
destra.

\begin{affianco}[0.55][c]
\begin{lstlisting}[style=c]
printf("abc\n");
printf("%c%c%c\n", 'a', 'b', 'c');
printf("%s\n", "abc");
printf("%c%3c%5c\n", 'A', 'B', 'C');
\end{lstlisting}
\risultato
\begin{console}
abc
abc
abc
A  B    C
\end{console}
\end{affianco}

Nel codice, \code{'a'} fra apici singoli è una \textbf{costante carattere}
(un solo carattere), mentre \code{"abc"} fra virgolette è una
\textbf{stringa}. I caratteri di conversione principali di \code{printf}
sono:

\begin{center}
\renewcommand{\arraystretch}{1.3}
\begin{tabular}{@{}>{\FiraCodeNL}c l l@{}}
  \toprule
  \normalfont\textbf{conversione} & \textbf{stampa} & \textbf{esempio} \\
  \midrule
  \%c & un carattere & \code{printf("\%c", 'A')} $\to$ \code{A} \\
  \%d & un intero, in decimale & \code{printf("\%d", 42)} $\to$ \code{42} \\
  \%e & un reale, in notazione scientifica & \code{printf("\%e", 314159.0)} $\to$ \code{3.141590e+05} \\
  \%f & un reale & \code{printf("\%f", 3.14159)} $\to$ \code{3.141590} \\
  \%g & un reale, nel più corto fra \code{\%e} e \code{\%f} & \code{printf("\%g", 314159.0)} $\to$ \code{314159} \\
  \%s & una stringa & \code{printf("\%s", "abc")} $\to$ \code{abc} \\
  \bottomrule
\end{tabular}
\end{center}

Fra \code{\%} e il carattere di conversione si possono indicare l'ampiezza
del campo e, per i reali, il numero di cifre decimali (la
\textit{precisione}); un segno meno allinea a sinistra. Le parentesi quadre
nell'esempio servono solo a mostrare dove comincia e finisce il campo:

\begin{affianco}[0.55][c]
\begin{lstlisting}[style=c]
printf("[%d] [%5d] [%-5d]\n", 42, 42, 42);
printf("[%f] [%.2f] [%10.3f]\n", 3.14159, 3.14159, 3.14159);
printf("[%e] [%g] [%g]\n", 314159.0, 314159.0, 0.0000314);
\end{lstlisting}
\risultato
\begin{console}[Terminale]
[42] [   42] [42   ]
[3.141590] [3.14] [     3.142]
[3.141590e+05] [314159] [3.14e-05]
\end{console}
\end{affianco}

Si noti che \code{\%f} stampa sempre 6 cifre decimali, se non si indica
altro, e che la precisione \textbf{arrotonda}: 3.14159 con 3 decimali
diventa 3.142.

\subsection{\code{scanf} e l'operatore \code{\&}}

\code{scanf} è l'analoga di \code{printf} per l'input. Il primo parametro è
la stringa di controllo; l'input è sempre una sequenza di caratteri (ciò che
si digita sulla tastiera), e le conversioni specificano come i caratteri
devono essere interpretati. I parametri successivi, però, non sono variabili
ma \textbf{indirizzi}: gli indirizzi delle variabili in cui scrivere i
valori letti.

\begin{lstlisting}[style=c]
scanf("%d", &x);
\end{lstlisting}

\code{\&} è l'\term{operatore di indirizzo}: \code{\&x} è l'indirizzo della
variabile \code{x}, il “numero civico” della sua cella di memoria
(figura~\ref{fig:scanf-indirizzo}). Perché serve? Perché \code{scanf} deve
\textbf{scrivere} nella variabile, e per farlo deve sapere dove si trova.
Scrivendo \code{scanf("\%d", x)} le si passerebbe il \textit{valore} di
\code{x} (un numero qualunque, ancora indefinito), e \code{scanf} non
saprebbe dove mettere il dato letto; il motivo profondo lo vedremo alla
fine del capitolo, con il passaggio dei parametri per valore.

\begin{figure}[H]
  \centering
  \begin{tikzpicture}[font=\Lato\scriptsize]
    \node[fill=vscbg, rounded corners=3pt, inner sep=6pt, font=\FiraCodeNL\small] (cod) at (0,0)
      {\textcolor{vscfunc}{scanf}\textcolor{vscfg}{(}\textcolor{vscstring}{"\%d"}\textcolor{vscfg}{, }\textcolor{yellow-gradient-2}{\bfseries \&x}\textcolor{vscfg}{);}};
    \node[wtnode, font=\Lato\scriptsize, text width=3.4cm] (sc) at (6.0,0) {\textbf{scanf}\\legge i caratteri digitati, li converte in un intero e lo scrive all'indirizzo ricevuto};
    \memoria{m}{(11.8,0.72)}{996/\ldots/, 1000/42/x, 1004/\ldots/}
    \draw[wtarrow] (cod.east) -- node[wtlab, above=2pt, fill=none] {riceve \code{\&x}, cioè 1000} (sc.west);
    \draw[wtarrowacc] (sc.east) -- node[wtlab, above=2pt, fill=none] {scrive 42} ([xshift=-30pt]m-2.west);
    \node[mnota, anchor=north, align=center] at ([yshift=-4pt]m-3.south) {la variabile \code{x}\\all'indirizzo 1000};
  \end{tikzpicture}
  \caption{\code{scanf} riceve l'indirizzo di \code{x} e scrive lì il
           valore letto (l'indirizzo 1000 è solo un esempio).}
  \label{fig:scanf-indirizzo}
\end{figure}

Le conversioni principali di \code{scanf}:

\begin{center}
\renewcommand{\arraystretch}{1.3}
\begin{tabular}{@{}>{\FiraCodeNL}c l@{}}
  \toprule
  \normalfont\textbf{conversione} & \textbf{legge e scrive in una variabile di tipo} \\
  \midrule
  \%c  & \code{char} (un carattere) \\
  \%d  & \code{int} (un intero decimale) \\
  \%f  & \code{float} \\
  \%lf & \code{double} \\
  \%s  & una stringa (capitolo~\ref{chap:stringhe}) \\
  \bottomrule
\end{tabular}
\end{center}

\warning{\code{\%f} o \code{\%lf}?}{
  Con \code{printf}, \code{\%f} va bene sia per i \code{float} sia per i
  \code{double}. Con \code{scanf} no: per un \code{double} serve
  \code{\%lf}. Il motivo: \code{scanf} riceve solo un indirizzo, e dalla
  conversione deve capire quanti byte scrivere a quell'indirizzo (4 per un
  \code{float}, 8 per un \code{double}). Con la conversione sbagliata
  scrive un numero di byte sbagliato, e il valore letto è senza senso.}

\subsection{Un esempio: echo}

Il programma seguente legge tre caratteri, un intero, un \code{float} e un
\code{double}, e li ristampa (fa da “eco”):

\begin{lstlisting}[style=c, name=echo.c]
#include <stdio.h>

int main(void)
{
    char     c1, c2, c3;
    int      i;
    float    x;
    double   y;

    printf("\n%s\n%s", "Input three characters ", "an int, a float, and a double \n");
    scanf("%c%c%c%d%f%lf", &c1, &c2, &c3, &i, &x, &y);
    printf("\nHere is the data that you typed in:\n");
    printf("%3c%3c%3c\n%17d%17f%17f\n\n", c1, c2, c3, i, x, y);
    return 0;
}
\end{lstlisting}

Proviamolo due volte, con input quasi uguali: nel secondo c'è solo uno
spazio fra \code{B} e \code{C}.

\begin{console}
$ ./echo

Input three characters
an int, a float, and a double
<<ABC 3 55 77.7>>

Here is the data that you typed in:
  A  B  C
                3        55.000000        77.700000

\end{console}

\begin{console}
$ ./echo

Input three characters
an int, a float, and a double
<<AB C 3 55 77.7>>

Here is the data that you typed in:
  A  B
       1386527472         0.000000         0.000000

\end{console}

Il secondo risultato è disastroso, e il motivo è una regola di
\code{scanf}: quando legge \textbf{numeri} salta gli spazi che li
precedono, ma quando legge \textbf{caratteri} no, perché lo spazio
\textit{è} un carattere. La figura~\ref{fig:scanf-buffer} mostra quali
caratteri “mangia” ciascuna conversione.

\begin{figure}[H]
  \centering
  \begin{tikzpicture}[font=\FiraCodeNL\scriptsize,
      ch/.style={mcella, minimum width=0.46cm, minimum height=0.5cm, inner xsep=1pt},
      conv/.style={font=\FiraCodeNL\scriptsize\bfseries, text=accent!70!black},
      tit/.style={font=\Lato\scriptsize\bfseries, anchor=east}]
    \noindici
    \node[tit] at (-0.2,0) {input 1};
    \celle[ch]{a}{(0,0)}{A,B,C,\char"2423,3,\char"2423,5,5,\char"2423,7,7,.,7,\textbackslash n}
    \foreach \ia/\ib/\conv in {0/0/\%c, 1/1/\%c, 2/2/\%c, 3/4/\%d, 5/7/\%f, 8/12/\%lf} {
      \draw[wtbrace] ([yshift=-3pt]a-\ia.south west) -- ([yshift=-3pt]a-\ib.south east);
      \node[conv, anchor=north] at ([yshift=-0.32cm]$(a-\ia.south)!0.5!(a-\ib.south)$) {\conv};
    }
    \node[mnota, anchor=west] at ([xshift=0.4cm]a-13.east) {6 conversioni riuscite: \textbf{A}, \textbf{B}, \textbf{C}, \textbf{3}, \textbf{55}, \textbf{77.7}};
    \node[tit] at (-0.2,-2.0) {input 2};
    \celle[ch]{b}{(0,-2.0)}{A,B,\char"2423,C,\char"2423,3,\char"2423,5,5,\char"2423,7,7,.,7}
    \foreach \ia/\ib/\conv in {0/0/\%c, 1/1/\%c, 2/2/\%c} {
      \draw[wtbrace] ([yshift=-3pt]b-\ia.south west) -- ([yshift=-3pt]b-\ib.south east);
      \node[conv, anchor=north] at ([yshift=-0.32cm]$(b-\ia.south)!0.5!(b-\ib.south)$) {\conv};
    }
    \draw[wtbrace, draw=merr] ([yshift=-3pt]b-3.south west) -- ([yshift=-3pt]b-3.south east);
    \node[conv, anchor=north, text=merr] at ([yshift=-0.32cm]b-3.south) {\%d \merr};
    \node[mnota, anchor=west, text width=6.4cm] at ([xshift=0.4cm]b-13.east) {la terza \code{\%c} legge lo \textbf{spazio}; poi \code{\%d} trova
      una \code{C}, che non è un numero, e \code{scanf} si ferma: \code{i}, \code{x} e \code{y}
      restano con i valori a caso che avevano};
  \end{tikzpicture}
  \caption{I caratteri dell'input (\code{\char"2423} indica uno spazio) e la
           conversione che li consuma.}
  \label{fig:scanf-buffer}
\end{figure}

Nel secondo caso \code{i}, \code{x} e \code{y} non ricevono nulla, e
\code{printf} stampa ciò che contenevano: valori casuali, diversi da un
computer all'altro (e da un'esecuzione all'altra). Come accorgersene dal
programma? \code{scanf} \textbf{restituisce il numero di conversioni
riuscite}. Salvandolo in una variabile e stampandolo (è il programma
\code{echo\_conversions.c} delle lezioni) si ottiene 6 nel primo caso e 3 nel
secondo:

\begin{affianco}[0.55][c]
\begin{lstlisting}[style=c]
int conversions;
...
conversions = scanf("%c%c%c%d%f%lf", &c1, &c2, &c3, &i, &x, &y);
...
printf("Correct conversions. %d\n", conversions);
\end{lstlisting}
\risultato
\begin{console}
Correct conversions. 6
...
Correct conversions. 3
\end{console}
\end{affianco}

Controllare il valore restituito da \code{scanf} è il modo corretto di
verificare che l'input sia stato letto: lo useremo fra poco.

\section{Il controllo del flusso}

\subsection{Vero e falso in C}

Le istruzioni sono normalmente eseguite in sequenza, ma la maggior parte
degli algoritmi richiede di alterare il flusso sequenziale
(capitolo~\ref{chap:strutturata}). Il C offre \code{if} e
\code{if}-\code{else} per scegliere fra azioni alternative, \code{while} e
\code{for} per ripetere. Tutti richiedono di valutare un'espressione
\textbf{logica}, cioè vera o falsa, e in C la regola è semplicissima:

\begin{itemize}
  \item il \textbf{falso} è rappresentato dal valore \textbf{0};
  \item \textbf{ogni valore diverso da zero} è interpretato come
        \textbf{vero}.
\end{itemize}

\subsection{\code{if}}

La forma generale è

\begin{lstlisting}[style=c]
if (espressione) {
    istruzioni_if;
}
istruzione_successiva;
\end{lstlisting}

Se \code{espressione} è diversa da zero (vera) si eseguono
\code{istruzioni\_if} e poi \code{istruzione\_successiva}; se vale 0 (falsa)
si salta direttamente a \code{istruzione\_successiva}. È la selezione a una
via del capitolo precedente. Un esempio:

\begin{affianco}[0.4][c]
\begin{lstlisting}[style=c]
a = 1;
if (b == 3) {
    a = 5;
}
printf("%d", a);
\end{lstlisting}
\risultato
Il simbolo \code{==} è l'operatore di \textbf{uguaglianza}: verifica se
\code{b} è uguale a 3. Se è vero, si assegna 5 ad \code{a} e viene stampato
\textbf{5}; se è falso si esegue subito la stampa, e viene stampato
\textbf{1}.
\end{affianco}

\error{\code{==} e \code{=} sono due cose diverse}{
  \code{==} \textbf{confronta}, \code{=} \textbf{assegna}. Scrivere
  \code{if (b = 3)} non è un errore per il compilatore: assegna 3 a
  \code{b}, e siccome il valore dell'assegnazione è 3 (diverso da zero, cioè
  vero), il ramo viene eseguito \textbf{sempre}. È uno degli errori più
  frequenti in C; il capitolo~\ref{chap:flusso} ci torna in dettaglio.
  Compilando con \code{-Wall}, gcc avvisa.}

Fra le graffe può esserci una sequenza di istruzioni; con una sola
istruzione le graffe si possono omettere, ma conviene \textbf{metterle
sempre}: se in seguito si aggiunge un'istruzione al blocco, senza graffe
resterebbe fuori dall'\code{if}, con un errore difficile da vedere.

\subsection{\code{if}-\code{else}}

\begin{affianco}[0.4][c]
\begin{lstlisting}[style=c]
if (espressione) {
    istruzioni_if;
} else {
    istruzioni_else;
}
\end{lstlisting}
\risultato
Se \code{espressione} è diversa da zero si esegue il primo blocco, se è zero
il secondo: è la selezione a due vie. Uno dei due blocchi viene sempre
eseguito, mai entrambi.
\end{affianco}

Nell'esempio seguente, se \code{cnt} vale 0 il programma stampa
$2 + 3 + 5 = 10$, altrimenti $-1 - 2 - 3 = -6$. A destra la variante con
\code{if (cnt)}: non c'è nessun confronto, la decisione si basa solo sul
valore di \code{cnt} (0 è falso, il resto vero), e i risultati si
invertono.

\begin{affianco}[0.5]
\begin{lstlisting}[style=c]
if (cnt == 0) {
    a = 2;
    b = 3;
    c = 5;
} else {
    a = -1;
    b = -2;
    c = -3;
}
printf("%d", a + b + c);
\end{lstlisting}
\risultato
\begin{lstlisting}[style=c]
if (cnt) {
    a = 2;
    b = 3;
    c = 5;
} else {
    a = -1;
    b = -2;
    c = -3;
}
printf("%d", a + b + c);
\end{lstlisting}
\end{affianco}

\begin{center}
\renewcommand{\arraystretch}{1.3}
\begin{tabular}{@{}lcc@{}}
  \toprule
  & \textbf{con \code{cnt == 0}} & \textbf{con \code{cnt}} \\
  \midrule
  se \code{cnt} vale 0           & stampa 10 & stampa $-6$ \\
  se \code{cnt} vale, per esempio, 4 & stampa $-6$ & stampa 10 \\
  \bottomrule
\end{tabular}
\end{center}

\subsection{\code{while}}

\begin{affianco}[0.42]
\begin{lstlisting}[style=c, name=sum.c]
#include <stdio.h>

int main(void)
{
    int   i = 1, sum = 0;

    while (i <= 5) {
        sum = sum + i;
        i = i + 1;
    }
    printf("sum = %d\n", sum);
    return 0;
}
\end{lstlisting}
\risultato
\begin{console}
$ ./sum
sum = 15
\end{console}
Il programma somma i numeri da 1 a 5. Nella dichiarazione,
\code{int i = 1, sum = 0;} le due variabili vengono anche
\textbf{inizializzate}. \code{<=} significa “minore o uguale”.
\end{affianco}

Il \code{while} verifica se \code{i <= 5}: se la condizione è vera esegue
le istruzioni del blocco fino alla \code{\}}, poi il controllo
\textbf{torna all'inizio} del ciclo per ripetere il test. L'iterazione
continua finché il test non fallisce. È il ciclo con la condizione in testa
del capitolo~\ref{chap:strutturata}. Simuliamolo:

\begin{center}
\renewcommand{\arraystretch}{1.3}
\begin{tabular}{@{}cccc@{}}
  \toprule
  \textbf{giro} & \textbf{\code{i <= 5}?} & \textbf{\code{sum} dopo il giro} & \textbf{\code{i} dopo il giro} \\
  \midrule
  1 & $1 \le 5$, vero & 0 + 1 = 1 & 2 \\
  2 & $2 \le 5$, vero & 1 + 2 = 3 & 3 \\
  3 & $3 \le 5$, vero & 3 + 3 = 6 & 4 \\
  4 & $4 \le 5$, vero & 6 + 4 = 10 & 5 \\
  5 & $5 \le 5$, vero & 10 + 5 = 15 & 6 \\
  6 & $6 \le 5$, \textbf{falso} & si esce: \code{sum} vale 15 & \\
  \bottomrule
\end{tabular}
\end{center}

La forma generale è

\begin{lstlisting}[style=c]
while (espressione) {
    istruzioni_while;
}
istruzioni;
\end{lstlisting}

Se \code{espressione} è diversa da zero si eseguono
\code{istruzioni\_while} e poi si verifica di nuovo \code{espressione}; il
processo continua finché \code{espressione} vale 0. Perché il ciclo
finisca, è \textbf{necessaria} un'istruzione dentro il ciclo che possa
cambiare il valore dell'espressione (qui \code{i = i + 1}): senza, il test
resterebbe vero per sempre.

\subsection{Assegnazioni abbreviate e incremento}

Istruzioni come \code{sum = sum + i;} sono così frequenti che il C ha una
forma abbreviata: \code{sum += i;}. In generale, un'istruzione della forma
\codel{variabile op= espressione;} (con \code{op} uno fra \code{+}, \code{-},
\code{*}, \code{/}) equivale a \codel{variabile = variabile op espressione;}.

Per aggiungere o togliere 1 ci sono gli \textbf{operatori di incremento e
decremento}:

\begin{center}
\renewcommand{\arraystretch}{1.3}
\begin{tabular}{@{}>{\FiraCodeNL}l l@{}}
  \toprule
  \normalfont\textbf{forma} & \textbf{come istruzione a sé, equivale a} \\
  \midrule
  ++i \textnormal{e} i++ & \code{i = i + 1;} \\
  --i \textnormal{e} i-- & \code{i = i - 1;} \\
  \bottomrule
\end{tabular}
\end{center}

Il ciclo di \code{sum.c} si scrive quindi anche così:

\begin{lstlisting}[style=c]
while (i <= 5) {
    sum += i;
    ++i;
}
\end{lstlisting}

\warning{\code{++} dentro le espressioni}{
  Come istruzioni isolate, \code{++i} e \code{i++} fanno la stessa cosa. Usati
  \textbf{dentro} un'espressione più grande (per esempio \code{a = i++;}) si
  comportano in modo diverso, e il loro funzionamento è più complicato: lo
  vedremo con calma nel capitolo~\ref{chap:sintassi}. Per ora li useremo
  solo da soli.}

\subsection{\code{for}}

Il ciclo \code{for} ha la forma

\begin{affianco}[0.48][c]
\begin{lstlisting}[style=c]
for (expr1; expr2; expr3) {
    istruzioni_for;
}
\end{lstlisting}
\risultato
\begin{lstlisting}[style=c]
expr1;
while (expr2) {
    istruzioni_for;
    expr3;
}
\end{lstlisting}
\end{affianco}

ed è equivalente al \code{while} a destra. Tipicamente \code{expr1} fa
un'assegnazione iniziale, \code{expr2} un test, \code{expr3} incrementa una
variabile: sono i tre passi del ciclo su contatore del
capitolo~\ref{chap:strutturata} (inizializzazione, test, modifica), riuniti
in un'unica riga. La somma da 1 a 5 diventa:

\begin{affianco}[0.48][c]
\begin{lstlisting}[style=c]
for (i = 1; i <= 5; i++) {
    sum += i;
}
\end{lstlisting}
\risultato
\begin{immagine}
\begin{tikzpicture}[font=\Lato\scriptsize]
  \node[fcproc, fill=green-gradient-1!40!white] (a) at (0,0) {\code{i = 1}};
  \node[fcgiunzione] (j) at (0,-0.6) {};
  \node[fcdec, fill=yellow-gradient-1!60!white] (t) at (0,-1.35) {\code{i <= 5}?};
  \node[fcproc] (b) at (0,-2.35) {\code{sum += i}};
  \node[fcproc, fill=accent!10!white] (c) at (0,-3.15) {\code{i++}};
  \draw[fcfreccia] (a) -- (j); \draw[fcfreccia] (j) -- (t);
  \draw[fcfreccia] (t) -- node[fcsi, right] {vero} (b);
  \draw[fcfreccia] (b) -- (c);
  \draw[fcfreccia, draw=accent!70!black] (c.south) -- (0,-3.6) -- (1.6,-3.6) -- (1.6,-0.6) -- (j.east);
  \draw[fcfreccia] (t.west) -- node[fcno, above] {falso} (-1.6,-1.35) -- (-1.6,-3.9);
\end{tikzpicture}
\end{immagine}
\end{affianco}

\section{Un esempio completo: \code{running\_sum}}

\subsection{Il programma}

Il programma seguente legge una sequenza di numeri e, per ciascuno, stampa
una riga di tabella con il conteggio, il numero letto, il minimo, il
massimo, la somma e la media dei numeri letti fino a quel momento.

\begin{lstlisting}[style=c, name=running_sum.c]
/* Compute the minimum, maximum, sum, and average. */

#include <stdio.h>
#include <stdlib.h>

int main(void)
{
    int      i;
    double   x, min, max, sum, avg;

    if (scanf("%lf", &x) != 1) {
        printf("No data found - bye!\n");
        exit(1);
    }
    min = max = sum = avg = x;
    printf("%5s%9s%9s%9s%12s%12s\n%5s%9s%9s%9s%12s%12s\n\n",
           "Count", "Item", "Min", "Max", "Sum", "Average",
           "-----", "----", "---", "---", "---", "-------");
    printf("%5d%9.1f%9.1f%9.1f%12.3f%12.3f\n",
           1, x, min, max, sum, avg);
    for (i = 2; scanf("%lf", &x) == 1; ++i) {
        if (x < min)
            min = x;
        else if (x > max)
            max = x;
        sum += x;
        avg = sum / i;
        printf("%5d%9.1f%9.1f%9.1f%12.3f%12.3f\n",
               i, x, min, max, sum, avg);
    }
    return 0;
}
\end{lstlisting}

\subsection{Riga per riga}

\begin{itemize}
  \item \code{if (scanf("\%lf", \&x) != 1)}: \code{scanf} restituisce il
        numero di conversioni riuscite; qui ne è richiesta una sola, quindi
        un valore diverso da 1 (\code{!=} significa “diverso”) vuol dire
        che non è stato letto niente. In quel caso il programma stampa un
        messaggio ed esce con \code{exit(1)}: \code{exit}, il cui prototipo
        sta in \code{stdlib.h}, termina il programma restituendo al sistema
        operativo il valore indicato (qui 1, “qualcosa è andato storto”).
  \item \code{min = max = sum = avg = x;}: un'assegnazione in C è anche
        un'espressione, che \textbf{vale} il valore assegnato. Le
        assegnazioni a catena si eseguono da destra a sinistra:
        \code{avg = x}, poi \code{sum = (avg = x)}, e così via. Tutte e
        quattro le variabili ricevono il primo numero letto.
  \item Le due \code{printf} con \code{\%5s}, \code{\%9s}\dots\ stampano
        l'intestazione della tabella: stringhe costanti in campi di
        larghezza fissa, così le colonne restano allineate.
  \item \code{for (i = 2; scanf("\%lf", \&x) == 1; ++i)}: un ciclo \code{for}
        in cui il \textbf{test è la lettura stessa}. Il ciclo continua finché
        \code{scanf} riesce a leggere un numero, cioè finché c'è input; non
        si sa in anticipo quante volte girerà. \code{i} conta i numeri letti
        e parte da 2 perché il primo è già stato letto prima del ciclo.
  \item \code{if (x < min) min = x; else if (x > max) max = x;}: aggiorna, se
        necessario, il minimo o il massimo. L'\code{else} contiene a sua
        volta un \code{if}. (Si poteva anche scrivere con due \code{if}
        indipendenti.)
  \item \code{avg = sum / i;}: qui la divisione è fra un \code{double} e un
        \code{int}, quindi è reale: nessuna trappola.
\end{itemize}

\subsection{Tre modi di dare l'input}

Il programma legge i numeri dallo \term{standard input}, il “canale” da cui
un programma riceve i dati. Normalmente è la tastiera, e allora i numeri si
possono scrivere uno alla volta, premendo \tasto{Invio} dopo ciascuno, oppure
tutti su una riga; ma come fa il programma a sapere che l'input è finito?
Bisogna segnalarlo con \tasti{Ctrl,D} all'inizio di una riga: \code{scanf}
non riesce più a leggere, restituisce un valore diverso da 1, e il ciclo
finisce.

\begin{console}
$ ./running_sum
<<3>>
Count     Item      Min      Max         Sum     Average
-----     ----      ---      ---         ---     -------

    1      3.0      3.0      3.0       3.000       3.000
<<-5 7>>
    2     -5.0     -5.0      3.0      -2.000      -1.000
    3      7.0     -5.0      7.0       5.000       1.667
<<Ctrl+D>>
$
\end{console}

(\tasti{Ctrl,D} in realtà non compare sullo schermo: qui è indicato per
chiarezza.) Il terzo modo è il più comodo per provare un programma: scrivere i dati in
un file e usarlo \textbf{al posto} della tastiera, con la
\term{redirezione} \code{<} (figura~\ref{fig:redirezione}). Il file
\code{running\_sum.in} delle lezioni contiene i dieci numeri
\code{3 -5 7 -9 11 -13 15 -17 19 -21}.

\begin{console}
$ ./running_sum < running_sum.in
Count     Item      Min      Max         Sum     Average
-----     ----      ---      ---         ---     -------

    1      3.0      3.0      3.0       3.000       3.000
    2     -5.0     -5.0      3.0      -2.000      -1.000
    3      7.0     -5.0      7.0       5.000       1.667
    4     -9.0     -9.0      7.0      -4.000      -1.000
    5     11.0     -9.0     11.0       7.000       1.400
    6    -13.0    -13.0     11.0      -6.000      -1.000
    7     15.0    -13.0     15.0       9.000       1.286
    8    -17.0    -17.0     15.0      -8.000      -1.000
    9     19.0    -17.0     19.0      11.000       1.222
   10    -21.0    -21.0     19.0     -10.000      -1.000
\end{console}

\begin{figure}[H]
  \centering
  \begin{tikzpicture}[font=\Lato\scriptsize,
      prog/.style={wtnode, font=\Lato\scriptsize\bfseries, minimum width=2.6cm, minimum height=0.9cm},
      dev/.style={wtnodeplain, minimum width=2.3cm, minimum height=0.8cm},
      can/.style={font=\FiraCodeNL\scriptsize, text=accent!70!black}]
    \node[dev] (k) at (0,0) {tastiera};
    \node[prog] (p) at (4.9,0) {running\_sum};
    \node[dev] (s) at (9.8,0) {schermo};
    \draw[wtarrow] (k) -- node[can, above] {standard input} (p);
    \draw[wtarrow] (p) -- node[can, above] {standard output} (s);
    \node[font=\Lato\scriptsize\bfseries, anchor=west] at (11.4,0) {normalmente};
    \node[dev, draw=orange-gradient-6!80!black, fill=orange-gradient-1!30!white] (f) at (0,-1.5) {\FiraCodeNL running\_sum.in};
    \node[prog] (p2) at (4.9,-1.5) {running\_sum};
    \node[dev] (s2) at (9.8,-1.5) {schermo};
    \draw[wtarrow] (f) -- node[can, above] {<} (p2);
    \draw[wtarrow] (p2) -- (s2);
    \node[font=\Lato\scriptsize\bfseries, anchor=west, align=left] at (11.4,-1.5) {\FiraCodeNL ./running\_sum < running\_sum.in};
    \node[dev] (k3) at (0,-3.0) {tastiera};
    \node[prog] (p3) at (4.9,-3.0) {running\_sum};
    \node[dev, draw=orange-gradient-6!80!black, fill=orange-gradient-1!30!white] (s3) at (9.8,-3.0) {\FiraCodeNL risultati.txt};
    \draw[wtarrow] (k3) -- (p3);
    \draw[wtarrow] (p3) -- node[can, above] {>} (s3);
    \node[font=\Lato\scriptsize\bfseries, anchor=west, align=left] at (11.4,-3.0) {\FiraCodeNL ./running\_sum > risultati.txt};
  \end{tikzpicture}
  \caption{La redirezione: con \code{<} un file prende il posto della
           tastiera, con \code{>} un file prende il posto dello schermo.}
  \label{fig:redirezione}
\end{figure}

È la shell a collegare il file al programma: il programma non se ne accorge
nemmeno, continua a leggere dallo standard input. Simmetricamente, con
\code{>} si può mandare l'output in un file invece che sullo schermo, e i due
si possono combinare: \code{./running\_sum < running\_sum.in > risultati.txt}.
Infine, se non c'è nessun dato, il programma esce con il suo messaggio, e il
valore restituito al sistema è 1:

\begin{console}
$ ./running_sum < /dev/null
No data found - bye!
$ echo $?
1
\end{console}

(\code{/dev/null} è un file speciale di Linux, sempre vuoto.)

\subsection{Una versione senza la lettura iniziale}

Nel programma si fa una prima lettura \textbf{prima} del ciclo, poi le altre
\textbf{nel} ciclo: lo stesso codice di lettura compare due volte. Quando è
possibile, è più elegante mettere tutto dentro il ciclo. Ecco la versione
delle lezioni con un \code{while}, in cui \code{cont} conta i numeri letti e
il caso “nessun dato” si controlla alla fine:

\begin{lstlisting}[style=c, name=running_sum2.c]
/* Compute the minimum, maximum, sum, and average. */

#include <stdio.h>

int main(void)
{
    int      cont = 0;
    double   x, min, max, sum, avg;

    while (scanf("%lf", &x) == 1) {
        cont = cont + 1;
        if (cont == 1) {
            sum = avg = min = max = x;
            printf("%5s%9s%9s%9s%12s%12s\n%5s%9s%9s%9s%12s%12s\n\n",
                   "Count", "Item", "Min", "Max", "Sum", "Average",
                   "-----", "----", "---", "---", "---", "-------");
        } else {
            if (x < min) min = x;
            else if (x > max) max = x;
            sum += x;
            avg = sum / cont;
        }
        printf("%5d%9.1f%9.1f%9.1f%12.3f%12.3f\n",
               cont, x, min, max, sum, avg);
    }
    if (cont == 0) {
        printf("No data found - bye!\n");
    }
    return 0;
}
\end{lstlisting}

L'output è identico. Di solito una variabile come \code{i} è la variabile di
un ciclo su contatore; qui \code{cont} conta gli elementi letti, e il primo
elemento è trattato a parte dentro il ciclo (\code{cont == 1}). Non serve
più \code{stdlib.h}, perché non si usa \code{exit}.

\section{Quando la compilazione fallisce}

I messaggi del compilatore sono importanti e vanno letti; a volte però
possono creare confusione. A lezione si compila il file
\code{running\_sum2Errore.c}, identico a quello appena visto tranne che per
un dettaglio: nella prima riga manca la barra iniziale del commento.

\begin{lstlisting}[style=c, name=running_sum2Errore.c]
* Compute the minimum, maximum, sum, and average. */

#include <stdio.h>
...
\end{lstlisting}

Il risultato è una lunga, deprimente lista di \textbf{51 errori}, molti dei
quali addirittura negli header file di sistema:

\begin{console}
$ gcc -Wall -o running_sum2 running_sum2Errore.c
running_sum2Errore.c:1:11: error: expected '=', ',', ';', 'asm' or '__attribute__' before 'the'
    1 | * Compute the minimum, maximum, sum, and average. */
      |           ^~~
running_sum2Errore.c:1:11: error: unknown type name 'the'
In file included from /usr/include/stdio.h:44,
                 from running_sum2Errore.c:3:
/usr/include/x86_64-linux-gnu/bits/types/struct_FILE.h:95:3: error: unknown type name 'size_t'
   95 |   size_t __pad5;
      |   ^~~~~~
...
\end{console}

Che cosa è successo? Senza \code{/*} la prima riga non è più un commento: il
compilatore la legge come codice C e cerca di darle un senso
(\code{* Compute the\dots}), fallendo. Ma non sa dove finisce l'errore, e
continua a leggere il resto del file “fuori fase”: anche il contenuto di
\code{stdio.h}, perfettamente corretto, viene letto nel contesto sbagliato e
produce errori a cascata. La regola del capitolo~\ref{chap:ambiente} vale
più che mai: si guarda \textbf{il primo} errore, che qui indica
correttamente la riga 1, e si ignorano gli altri finché non lo si è
corretto.

\info{Una curiosità}{
  Il file \code{running\_sum2.c} distribuito con le lezioni ha lo stesso
  difetto di \code{running\_sum2Errore.c} (i due file sono identici):
  per compilarlo bisogna aggiungere la \code{/} mancante all'inizio.}

\section{Le funzioni}

\subsection{Dichiarare prima di usare}

Una \term{funzione} è un pezzo di codice che rappresenta un blocco del
processo di soluzione del problema: è il sotto-programma del
capitolo~\ref{chap:strutturata}. Un programma C è costituito da una o più
funzioni, in uno o più file sorgente. Una di esse è necessariamente
\code{main}, da cui comincia l'esecuzione; le altre sono richiamate da
\code{main} o da altre funzioni.

In C le funzioni devono essere \textbf{dichiarate prima di usarle}.
Supponiamo di voler usare la funzione \code{pow} della libreria matematica:
la chiamata \code{pow(x, y)} restituisce $x^y$. La sua dichiarazione è

\begin{lstlisting}[style=c]
double pow(double x, double y);
\end{lstlisting}

Questo tipo di dichiarazione si chiama \term{prototipo} della funzione: dice
il tipo del valore restituito (\code{double}), il nome e il tipo dei
parametri. I nomi dei parametri, nel prototipo, sono ignorati dal
compilatore: \code{double pow(double, double);} è equivalente, perché al
compilatore interessano solo i tipi. È comunque buona norma metterli,
soprattutto se sono significativi.

La forma generale di un prototipo è
\code{tipo nome\_funzione(lista dei tipi dei parametri);}. La parola chiave
\code{void} si usa quando la funzione non ha parametri
(\code{int f(void);}) o non restituisce niente (\code{void f(int);}). Se
necessario, al momento della chiamata gli argomenti vengono convertiti al
tipo indicato nel prototipo: chiamando \code{pow(2, 10)} con due interi,
questi diventano \code{2.0} e \code{10.0}.

\subsection{L'esempio maxmin}

Il programma seguente stampa delle informazioni su se stesso, legge un intero
\code{n}, legge \code{n} numeri reali e ne trova il massimo e il minimo.

\begin{lstlisting}[style=c, name=maxmin.c]
#include <stdio.h>

float   maximum(float x, float y);
float   minimum(float x, float y);
void    prn_info(void);

int main(void)
{
    int     i, n;
    float   max, min, x;

    prn_info();
    printf("Input n:  ");
    scanf("%d", &n);
    printf("\nInput %d real numbers:  ", n);
    scanf("%f", &x);
    max = min = x;
    for (i = 2; i <= n; ++i) {
        scanf("%f", &x);
        max = maximum(max, x);
        min = minimum(min, x);
    }
    printf("\n%s%11.3f\n%s%11.3f\n\n",
           "Maximum value:", max,
           "Minimum value:", min);
    return 0;
}

float maximum(float x, float y)
{
    if (x > y)
        return x;
    else
        return y;
}

float minimum(float x, float y)
{
    if (x < y)
        return x;
    else
        return y;
}

void prn_info(void)
{
    printf("\n%s\n%s\n\n",
           "This program reads an integer value for n, and then",
           "processes n real numbers to find max and min values.");
}
\end{lstlisting}

\begin{console}
$ ./maxmin

This program reads an integer value for n, and then
processes n real numbers to find max and min values.

Input n:  <<5>>

Input 5 real numbers:  <<737.7799  -11.2e+3  -777  0.001  3.14159>>

Maximum value:    737.780
Minimum value: -11200.000

\end{console}

Si noti che \code{scanf} con \code{\%f} accetta anche la notazione
scientifica: \code{-11.2e+3} è $-11{,}2 \cdot 10^3 = -11\,200$.

\subsection{Come funziona}

\begin{itemize}
  \item In cima al file ci sono i \textbf{prototipi} di \code{maximum},
        \code{minimum} e \code{prn\_info}. I primi due dicono che le funzioni
        hanno due parametri di tipo \code{float} e restituiscono un
        \code{float}; il terzo che \code{prn\_info} non ha parametri e non
        restituisce alcun valore.
  \item \code{prn\_info();} chiama la funzione, che a sua volta chiama
        \code{printf} per stampare le informazioni sul programma e sul suo
        uso.
  \item Il programma chiede e legge \code{n} (\code{\%d}: un intero
        decimale), poi legge il primo numero reale in \code{x}, e
        \code{max = min = x;} inizializza massimo e minimo con quel valore.
  \item A ogni giro del ciclo si legge un nuovo valore di \code{x};
        \code{max = maximum(max, x);} calcola il massimo fra \code{max} e
        \code{x} e lo assegna a \code{max}, e lo stesso per il minimo.
\end{itemize}

La \textbf{definizione} di \code{maximum} specifica che cosa la funzione fa
quando viene chiamata. È fatta di una \textbf{testa}, la parte prima della
prima graffa (\code{float maximum(float x, float y)}: il primo \code{float}
è il tipo del valore restituito, gli altri due i tipi dei parametri), e di un
\textbf{corpo}, la parte fra le graffe.

L'istruzione \code{return} \textbf{interrompe} la funzione e restituisce al
chiamante il valore della sua espressione; la forma generale è
\code{return espressione;}. In una funzione \code{void} si può scrivere
\code{return;} senza espressione, o lasciare che la funzione arrivi da sola
alla graffa finale, come fa \code{prn\_info}.

\subsection{Due \code{x} diverse}

Le funzioni \code{maximum} e \code{minimum} hanno entrambe i parametri
\code{x} e \code{y}, e una variabile \code{x} c'è anche in \code{main}. Non
c'è nessuna relazione fra loro: sono variabili \textbf{logicamente e
fisicamente diverse}, in celle di memoria diverse. Usare lo stesso nome
informa solo chi legge che il loro uso è simile. La
figura~\ref{fig:due-x} mostra la memoria durante la chiamata
\code{maximum(max, x)} al secondo giro.

\begin{figure}[H]
  \centering
  \begin{tikzpicture}[font=\Lato\scriptsize]
    \variabile{mi}{(1.4,0)}{i}{2}
    \variabile{mn}{(1.4,-0.6)}{n}{5}
    \variabile[minimum width=1.7cm]{mmax}{(1.55,-1.2)}{max}{737.7799}
    \variabile[minimum width=1.7cm]{mmin}{(1.55,-1.8)}{min}{737.7799}
    \variabile[minimum width=1.7cm]{mx}{(1.55,-2.4)}{x}{-11200}
    \scopo{2}{main}{main}{(mi-nome) (mn) (mmax) (mmin) (mx) (mx-nome)}
    \variabile[minimum width=1.7cm]{fx}{(8.0,-1.2)}{x}{737.7799}
    \variabile[minimum width=1.7cm]{fy}{(8.0,-1.8)}{y}{-11200}
    \scopo{3}{maximum}{maximum}{(fx-nome) (fx) (fy) (fy-nome)}
    \draw[mcerca] (mmax.east) -- node[wtlab, above, fill=none] {copia} (fx.west -| fx-nome.west);
    \draw[mcerca] (mx.east) -- node[wtlab, below, fill=none] {copia} (fy.west -| fy-nome.west);
  \end{tikzpicture}
  \caption{La \code{x} di \code{main} e la \code{x} di \code{maximum} sono
           due variabili distinte: \code{maximum} riceve una
           \textbf{copia} dei valori.}
  \label{fig:due-x}
\end{figure}

\subsection{Che cosa succede senza prototipi}

Le slide mostrano anche una versione di \code{maxmin} senza i prototipi in
cima al file. Con gcc il risultato è questo:

\begin{console}
$ gcc -Wall -o maxmin maxmin.c
maxmin.c:9:4: warning: implicit declaration of function (*@‘@*)prn_info(*@’@*)
maxmin.c:17:13: warning: implicit declaration of function (*@‘@*)maximum(*@’@*)
maxmin.c:18:13: warning: implicit declaration of function (*@‘@*)minimum(*@’@*)
maxmin.c:26:7: error: conflicting types for (*@‘@*)maximum(*@’@*); have (*@‘@*)float(float,  float)(*@’@*)
maxmin.c:34:7: error: conflicting types for (*@‘@*)minimum(*@’@*); have (*@‘@*)float(float,  float)(*@’@*)
\end{console}

Il compilatore legge il file dall'alto in basso. Quando incontra
\code{maximum(max, x)} non sa ancora nulla di \code{maximum}: per una vecchia
regola del C, “immagina” che sia una funzione che restituisce un
\code{int} (è la \textit{dichiarazione implicita} di cui avvisa il
warning). Più avanti trova la definizione vera, che restituisce un
\code{float}, e la contraddizione diventa un errore. Le versioni più recenti
di gcc (dalla 14) non accettano più le dichiarazioni implicite nemmeno come
avviso: la chiamata di una funzione non dichiarata è direttamente un errore.
La lezione è semplice: \textbf{ogni funzione va dichiarata prima dell'uso}.

Un modo alternativo, che va bene per i file piccoli, è mettere le
\textbf{definizioni} delle funzioni \textbf{prima} di \code{main}: quando il
compilatore arriva alla chiamata ha già letto la definizione, e la testa
della funzione fa da prototipo.

\subsection{Un punto e virgola di troppo}

Nelle prime slide su \code{maxmin} la funzione \code{maximum} contiene un
errore di battitura istruttivo:

\begin{affianco}[0.4]
\begin{lstlisting}[style=c]
float maximum(float x, float y)
{
    if (x > y) ;{
        return x;
    }
    else {
        return y;
    }
}
\end{lstlisting}
\risultato
\begin{console}
$ gcc -Wall -o maxmin maxmin.c
maxmin.c: In function (*@‘@*)maximum(*@’@*):
maxmin.c:31:4: warning: this (*@‘@*)if(*@’@*) clause does not guard... [-Wmisleading-indentation]
   31 |    if (x > y) ;{
      |    ^~
maxmin.c:34:4: error: (*@‘@*)else(*@’@*) without a previous (*@‘@*)if(*@’@*)
   34 |    else {
      |    ^~~~
\end{console}
\end{affianco}

Il \code{;} dopo la condizione è un'\textbf{istruzione vuota}: l'\code{if}
controlla quella, cioè nulla. Il blocco \code{\{ return x; \}} che segue non
fa più parte dell'\code{if}: è un blocco qualunque, eseguito sempre. E
quando arriva l'\code{else}, l'\code{if} a cui dovrebbe appartenere è già
concluso: “\code{else} senza un \code{if} precedente”. gcc avvisa anche
che l'indentazione è \textit{ingannevole}: il codice sembra dire una cosa e
ne fa un'altra.

\section{Il passaggio dei parametri per valore}

Esistono due modi di passare i parametri a una funzione:

\begin{itemize}
  \item il \textbf{passaggio per valore}: viene creata una \textbf{copia}
        della variabile passata come parametro;
  \item il \textbf{passaggio per riferimento}: viene passata la variabile
        stessa.
\end{itemize}

\textbf{In C i parametri sono passati solo per valore.} Il programma
seguente lo dimostra: la funzione prova a cambiare il valore di \code{a},
ma \code{main} non se ne accorge.

\begin{affianco}[0.5]
\begin{lstlisting}[style=c, name=try_to.c]
#include <stdio.h>

int main(void)
{
    int    a = 1;
    void   try_to_change_it(int);

    printf("%d\n", a);     /* 1 is printed */
    try_to_change_it(a);
    printf("%d\n", a);     /* 1 is printed again! */
    return 0;
}

void try_to_change_it(int a)
{
    a = 777;
}
\end{lstlisting}
\risultato
\begin{console}
$ ./try_to
1
1
\end{console}
\wtnota{Il prototipo è scritto dentro \code{main}: è lecito, e vale solo
lì. Di solito lo si mette in cima al file.}
\end{affianco}

\begin{figure}[H]
  \centering
  \begin{tikzpicture}[font=\Lato\scriptsize]
    \foreach \col/\testo in {0/{1. prima della chiamata}, 1/{2. dentro la funzione, dopo \code{a = 777}}, 2/{3. dopo il ritorno}} {
      \node[font=\Lato\scriptsize\bfseries] at (\col*5.6+1.6,1.25) {\testo};
    }
    % 1
    \variabile{a1}{(1.6,0)}{a}{1}
    \scopo{2}{m1}{main}{(a1-nome) (a1)}
    % 2
    \begin{scope}[xshift=5.6cm]
      \variabile{a2}{(1.6,0)}{a}{1}
      \scopo{2}{m2}{main}{(a2-nome) (a2)}
      \variabile[fill=accent!10!white]{b2}{(1.6,-1.55)}{a}{777}
      \scopo{3}{f2}{try\_to\_change\_it}{(b2-nome) (b2)}
      \draw[mcerca] (a2.east) to[out=-20, in=20] node[wtlab, right, fill=none] {copia di 1} (b2.east);
    \end{scope}
    % 3
    \begin{scope}[xshift=11.2cm]
      \variabile{a3}{(1.6,0)}{a}{1}
      \scopo{2}{m3}{main}{(a3-nome) (a3)}
      \node[mnota, align=center] at (1.3,-1.55) {la copia non esiste più:\\\code{a} di \code{main} vale ancora 1};
    \end{scope}
  \end{tikzpicture}
  \caption{Il passaggio per valore: la funzione lavora su una copia, che
           scompare quando la funzione termina.}
  \label{fig:per-valore}
\end{figure}

Alla chiamata, il parametro \code{a} di \code{try\_to\_change\_it} (una
variabile diversa da quella di \code{main}, anche se ha lo stesso nome)
riceve una \textbf{copia} del valore 1. L'assegnazione \code{a = 777}
modifica la copia; quando la funzione termina la copia scompare, e
\code{main} ritrova la sua \code{a} intatta (figura~\ref{fig:per-valore}).

Ecco perché \code{scanf} ha bisogno di \code{\&x}: se ricevesse \code{x},
riceverebbe solo una copia del suo valore e non potrebbe modificare la
variabile. Ricevendo l'\textbf{indirizzo}, sa dove andare a scrivere. Come
una funzione possa usare un indirizzo per modificare una variabile del
chiamante (i \textit{puntatori}) è l'argomento del
capitolo~\ref{chap:puntatori}.

\insintesi{
  Un programma C è fatto di funzioni, e l'esecuzione comincia da
  \code{main}. Le righe che cominciano con \code{\#} sono per il
  preprocessore: \code{\#include} copia un header (che \textbf{descrive} le
  funzioni; il loro codice arriva dalla libreria al momento del linking),
  \code{\#define} definisce costanti simboliche. Le variabili si dichiarano
  con un tipo prima di usarle; \code{=} assegna (si legge da destra a
  sinistra); la divisione fra interi tronca. \code{printf} e \code{scanf}
  usano una stringa di controllo con le conversioni (\code{\%d}, \code{\%f},
  \code{\%lf}\dots); \code{scanf} vuole gli \textbf{indirizzi} delle
  variabili e restituisce il numero di conversioni riuscite. In C 0 è falso
  e ogni altro valore è vero; \code{if}, \code{while} e \code{for}
  realizzano selezione e iterazione. Le funzioni vanno dichiarate con un
  prototipo prima dell'uso, e ricevono sempre una \textbf{copia} dei
  parametri.}

\section{Riferimenti}

\begin{itemize}
  \item \textbf{A. Kelley, I. Pohl}, \textit{A Book on C}, Addison-Wesley:
        capitolo~1, da cui provengono gli esempi di questo capitolo.
  \item \textbf{Deitel \& Deitel}, \textit{Il linguaggio C}: capitoli 2--5.
  \item \textbf{printf} e \textbf{scanf} --- cppreference.com. \\
        \urlx{https://en.cppreference.com/w/c/io/fprintf} \\
        \urlx{https://en.cppreference.com/w/c/io/fscanf}
\end{itemize}
